\section*{अध्याय \ref{ch:b1:nernst} --- अपचयोपचय साम्य और नेर्न्स्ट समीकरण}

\begin{solution}{exo:b1:nernst:1}
\ce{NH3} में N: $-$III; \ce{NO3-} में: $+$V। \ce{SO4^2-} में S: $+$VI।
\ce{H2O2} में O: $-$I। \ce{Cr2O7^2-} में Cr: $+$VI। \ce{ClO-} में Cl: $+$I।
\ce{Fe3O4} में Fe: औसतन $+8/3$ (एक आयरन(II) और दो आयरन(III))।
\ce{CH3OH} में C: $x + 4(+1) + (-2) = 0$, $-$II।
\end{solution}

\begin{solution}{exo:b1:nernst:2}
\ce{Cr2O7^2- + 14H+ + 6e- -> 2Cr^3+ + 7H2O};
\ce{H2O2 + 2H+ + 2e- -> 2H2O};
\ce{O2 + 2H+ + 2e- -> H2O2};
\ce{SO4^2- + 4H+ + 2e- -> SO2 + 2H2O}।
\end{solution}

\begin{solution}{exo:b1:nernst:3}
$E = -0.76 + 0.030\log[\ce{Zn^2+}]$; $E = 1.36 + 0.030\log(p(\ce{Cl2})/[\ce{Cl-}]^2)$;
$E = 1.23 + 0.015\log(p(\ce{O2})[\ce{H+}]^4)$; $E = 0.27 - 0.059\log[\ce{Cl-}]$
(दाब बार में)। \qty{0.010}{mol/L} ज़िंक सल्फ़ेट में:
$E = -0.76 + 0.030 \times (-2) = -0.82$~V।
\end{solution}

\begin{solution}{exo:b1:nernst:4}
$E(\ce{Zn}) = -0.76 + 0.030\log 0.10 = -0.79$~V, $E(\ce{Cu}) = 0.34 +
0.030\log 0.010 = 0.28$~V। वोल्टता $0.28 - (-0.79) = 1.07$~V। ज़िंक, जो
निम्नतर विभव पर है, ऋण ध्रुव है, जहाँ ऑक्सीकरण होता है: यही
\omterm{def:b1:nernst:cell}{ऐनोड} है।
\end{solution}

\begin{solution}{exo:b1:nernst:5}
\ce{MnO4- + 5Fe^2+ + 8H+ -> Mn^2+ + 5Fe^3+ + 4H2O}, $n = 5$,
$\log K = 5(1.51 - 0.77)/0.059 = 63$: \omterm{def:b1:extent-q-and-k:quantitative}{मात्रात्मक} और तीव्र, एक प्रसिद्ध
अनुमापन (परमैंगनेट स्वयं अपना सूचक है)।
\end{solution}

\begin{solution}{exo:b1:nernst:6}
जिन धातुओं का $E^\circ$ 0 से नीचे है वे अभिक्रिया करती हैं: Mg, Zn, Fe, Ni, Pb (लेड केवल
धीरे-धीरे: अंतर छोटा है); Cu और Ag नहीं। \ce{Mg + 2H+ -> Mg^2+ + H2}, $\log K = 2 \times 2.36/0.059 = 80$;
\ce{Fe + 2H+ -> Fe^2+ + H2}, $\log K = 2 \times 0.41/0.059 = 13.9$।
\end{solution}

\begin{solution}{exo:b1:nernst:7}
$n = 2$, $\log K = 2(0.53 - 0.02)/0.059 = 17$: \omterm{def:b1:extent-q-and-k:quantitative}{मात्रात्मक}।
\end{solution}

\begin{solution}{exo:b1:nernst:8}
$3E^\circ(\ce{Fe^3+}/\ce{Fe}) = 2(-0.41) + 0.77$, $E^\circ = -0.02$~V।
\ce{2Fe^3+ + Fe -> 3Fe^2+}: \omterm{def:b1:nernst:couple}{ऑक्सीकारक} \ce{Fe^3+} (0.77) \omterm{def:b1:nernst:couple}{अपचायक}
\ce{Fe} ($-0.41$) से ऊपर है, $n = 2$, $\log K = 2(0.77 + 0.41)/0.059 = 40$। आयरन
का बुरादा आयरन(II) के विलयन को आयरन(III) में बदलने से रोकता है।
\end{solution}

\begin{solution}{exo:b1:nernst:9}
\ce{O2 + 4H+ + 4e- -> 2H2O}: $E = 1.23 + \frac{0.059}{4}\log [\ce{H+}]^4 =
1.23 - 0.059\,\mathrm{pH}$। \ce{2H+ + 2e- -> H2}: $E = \frac{0.059}{2}\log
[\ce{H+}]^2 = -0.059\,\mathrm{pH}$। हर pH पर अंतर $1.23$~V है:
ईंधन सेल की वोल्टता pH पर निर्भर नहीं करती।
\end{solution}

\begin{solution}{exo:b1:nernst:10}
$E^\circ(\ce{Cu+}/\ce{Cu}) = 0.52 > E^\circ(\ce{Cu^2+}/\ce{Cu+}) = 0.16$: ऊपरी
युग्म का \omterm{def:b1:nernst:couple}{ऑक्सीकारक} \ce{Cu+} निचले युग्म के \omterm{def:b1:nernst:couple}{अपचायक} \ce{Cu+} से
अभिक्रिया करता है। $\log K = (0.52 - 0.16)/0.059 = 6.1$, $K =
[\ce{Cu^2+}]/[\ce{Cu+}]^2$, इसलिए $[\ce{Cu+}] = \sqrt{0.10/10^{6.1}} =
\qty{2.8e-4}{mol/L}$।
\end{solution}

\begin{solution}{exo:b1:nernst:11}
$U = 0.059\log(0.10/\num{1.0e-3}) = 0.118$~V। धन ध्रुव सांद्र विलयन वाला
सिल्वर है, जहाँ \ce{Ag+} अपचित होता है; तनु विलयन में
सिल्वर ऑक्सीकृत होता है। सांद्र विलयन तनु होता जाता है
और तनु विलयन सांद्र; वोल्टता घटती है और दोनों सांद्रताएँ बराबर होने पर
शून्य हो जाती है।
\end{solution}

\begin{solution}{exo:b1:nernst:12}
\ce{H2O2}, \ce{H2O2}/\ce{H2O} (1.76~V) का \omterm{def:b1:nernst:couple}{ऑक्सीकारक} है और \ce{O2}/\ce{H2O2} (0.69~V)
का \omterm{def:b1:nernst:couple}{अपचायक}: वह स्वयं से अभिक्रिया करता है, \ce{2H2O2 -> 2H2O + O2},
$\log K = 2(1.76 - 0.69)/0.059 = 36$। अभिक्रिया अनुकूल है पर \omterm{def:b1:elementary-steps:catalyst}{उत्प्रेरक} के बिना
बहुत धीमी (\cref{ch:b1:elementary-steps}); साफ़ बोतलें
और ठंडा, अँधेरा भंडारण इसे धीमा बनाए रखते हैं।
\end{solution}

\begin{solution}{pb:b1:nernst:1}
\textbf{1.} Ag: 0, $+$I, $+$I; Cl: दोनों में $-$I।
\textbf{2.} \ce{AgCl(s) + e- -> Ag(s) + Cl-}।
\textbf{3.} $E = E^\circ(\ce{AgCl}/\ce{Ag}) - 0.059\log[\ce{Cl-}]$।
\textbf{4.} दोनों ठोसों की \omterm{def:b1:extent-q-and-k:activity}{सक्रियता} 1 है; विलयन में केवल क्लोराइड आयन
है।
\textbf{5.} $E^\circ = 0.22$~V।
\textbf{6.} $0.22 + 0.118 = 0.34$~V।
\textbf{7.} \ce{Pt} $|$ \ce{H2} (\qty{1}{bar}) $|$ \ce{H+} (\omterm{def:b1:extent-q-and-k:activity}{सक्रियता} 1)
$\|$ \ce{Cl-} (\qty{0.10}{mol/L}) $|$ \ce{AgCl} $|$ \ce{Ag}।
\textbf{8.} सिल्वर इलेक्ट्रोड: $0.22 + 0.059 = 0.28$~V $> 0$।
\textbf{9.} $0.28$~V।
\textbf{10.} \ce{2AgCl + H2 -> 2Ag + 2H+ + 2Cl-}।
\textbf{11.} $\log K = 2 \times 0.22/0.059 = 7.5$।
\textbf{12.} प्लैटिनम (ऋण ध्रुव, जहाँ \ce{H2} ऑक्सीकृत होता है) से
सिल्वर की ओर।
\textbf{13.} $0.27 - 0.059\log 1.0 = 0.27$~V।
\textbf{14.} $E(\ce{Ag}) = 0.22 + 0.177 = 0.40$~V $> 0.27$: सिल्वर
इलेक्ट्रोड धनात्मक है, $U = 0.13$~V।
\textbf{15.} $U = 0.22 - 0.059\log[\ce{Cl-}] - 0.27 = -0.05 -
0.059\log[\ce{Cl-}]$ (वोल्ट में)।
\textbf{16.} $\log[\ce{Cl-}] = -(0.115 + 0.05)/0.059 = -2.80$,
$[\ce{Cl-}] = \qty{1.6e-3}{mol/L}$, अर्थात् \qty{57}{mg/L}।
\textbf{17.} $\Delta c/c = \ln 10 \times 0.001/0.059 = 0.039$: प्रति मिलीवोल्ट लगभग
\qty{4}{\percent}।
\textbf{18.} जब नमूने का क्लोराइड उतने के तुल्य हो जाए जितना
सिल्वर क्लोराइड स्वयं मुक्त करता है, $\sqrt{K_s} = \qty{1.3e-5}{mol/L}$;
इलेक्ट्रोड इससे काफ़ी ऊपर, लगभग \qty{1e-4}{mol/L} से, प्रयुक्त होता है।
\textbf{19.} $E = 0.80 + 0.059\log[\ce{Ag+}]$।
\textbf{20.} $[\ce{Ag+}] = K_s/[\ce{Cl-}]$।
\textbf{21.} $E = 0.80 + 0.059\log K_s - 0.059\log[\ce{Cl-}] = 0.80 -
0.059\,\mathrm{p}K_s - 0.059\log[\ce{Cl-}]$।
\textbf{22.} $E^\circ(\ce{AgCl}/\ce{Ag}) = E^\circ(\ce{Ag+}/\ce{Ag}) -
0.059\,\mathrm{p}K_s$ (\cref{prop:b1:nernst:second-kind})।
\textbf{23.} $0.80 - 0.059 \times 9.75 = 0.225$~V, सारणी के मान (0.22~V) से
सहमत, जो गिब्ज़ ऊर्जाओं से स्वतंत्र रूप से परिकलित हुआ था।
\textbf{24.} $E^\circ(\ce{AgCl}/\ce{Ag}) = $ \textbf{0.22~V}।
\end{solution}
